\chapter{Хэрэгжүүлэлт ба туршилт}
\section{Өгөгдлийн урьдчилсан боловсруулалт}
Цугласан дата дээр \texttt{fasttext} болон \texttt{pandas} сангуудыг ашиглан хэл таних болон шаардлагагүй үгс, тэмдэгтүүдийг цэвэрлэх пайплайныг Google Colab орчинд хэрэгжүүлсэн.

\section{Topic Modeling Хэрэгжүүлэлт}
BERTopic-ийн аргачлал дэмжин Mongolian-BERT sentence embedding-ийг UMAP болон HDBSCAN алгоритмуудтай хослуулан сургалаа. Урьдчилан тодорхойлсон ангиллын систем (Predefined topic classification) үүсгэн Cosine similarity ашигласнаар хурдан бөгөөд нарийвчлал өндөртэй үр дүн гарч байна. \texttt{sentence-transformers} сангийн тусламжтайгаар үгсийн вектор хоорондын зайг хэмжин кластеруудыг (сэдвүүдийг) тодорхойлсон.

\section{NER болон Network Graph туршилт}
\texttt{transformers} сангийн \texttt{Davlan/bert-base-multilingual-cased-ner-hrl} загварыг дуудан оруулж (inference pipeline), өгөгдсөн текстээс PERSON, ORG, LOC гэсэн ангиллаар токенуудыг ялгаж авав. Ялгаж авсан энэхүү өгөгдлөө ашиглан хэн, ямар байгууллагатай, хэзээ хамт дурдагдаж байгаа харилцааг Network Analysis буюу граф хэлбэрээр илэрхийлэх боломж бүрдлээ.

\section{Үнэлгээ болон дүгнэлт}
Сэдэв загварчлал болон NER-ийн нийлмэл алгоритмийг ашигласнаар зөвхөн үгсийн давтамжид тулгуурласан уламжлалт аргатай (Knowledge Tree Database) харьцуулахад контекстийг илүү сайн ойлгож, шинээр үүсч буй сэдвүүд болон нэршлийг өндөр нарийвчлалтайгаар ялгах боломжтойг хэрэгжүүлэлтийн шатны туршилтууд баталж байна.
